%══════════════════════════════════════════════════════════════════════
%  CHAPTER 3 — Related Work
%══════════════════════════════════════════════════════════════════════
\chapter{Related Work}
\label{ch:related}

This chapter reviews the literature on ML-based DDoS detection and
NFV orchestration, positions \padonap{} against eight recent systems
on a six-axis comparison matrix, and identifies three research gaps
that motivate the design choices made in this thesis.

%%─────────────────────────────────────────────────────────────────────
\section{ML-Based DDoS Detection Systems}
\label{sec:related-ml}

\textbf{TST-DDoS (IEEE Access 2025).}
TST-DDoS proposes a two-stage tree-based pipeline with engineered
flow features achieving an accuracy of 97.1\% and macro-F1 of 95.3\%
on CICDDoS2019. The system performs classification only and does not
integrate with a MANO or orchestration layer; it has no forecasting
or proactive capability.

\textbf{Lightweight ML (Scientific Reports 2025).}
This work evaluates a set of lightweight classifiers---Random Forest,
Decision Tree, and a shallow neural network---for DDoS detection in
resource-constrained edge nodes, achieving accuracy 96.2\% and F1
94.8\%. As with TST-DDoS, there is no orchestration integration,
multi-tenant SLA model, or proactive path.

\textbf{DAWN (Springer LNCS 2025).}
DAWN (DDoS-Aware NFV) introduces a hierarchical ML detector deployed
alongside a simplified NFV orchestrator. It partially addresses
orchestration (partial MANO) but does not implement a proactive forecast;
mitigation is initiated only after detection. Multi-tenant SLA is not
formalised.

%%─────────────────────────────────────────────────────────────────────
\section{Systems with Partial Orchestration or Forecasting}
\label{sec:related-orch}

\textbf{NetProbe (ICDCN 2025).}
NetProbe combines a Gradient Boosting classifier with a lightweight
probe-based orchestration mechanism. It achieves partial MANO
integration and includes a simple one-tier pre-positioning step,
but the forecasting horizon is a single step and no multi-tenant SLA
guarantee is provided. VNF boot latency is reported at approximately
3{,}000~ms, substantially higher than the 505~ms achieved by \padonap{}'s
pre-positioned rate-limiter.

\textbf{FlowGuard (NDSS 2022).}
FlowGuard introduces a single-horizon traffic forecaster integrated
with a flow-rule update mechanism. It partially addresses proactive
response but lacks a standards-compliant MANO interface and a multi-tenant
SLA model. The forecasting depth (single horizon, approximately 15~s)
is insufficient to cover full VNF boot cycles.

\textbf{Transformer+TCN (PLOS One 2025).}
This work combines a Transformer encoder with a Temporal Convolutional
Network (TCN) for multi-class DDoS classification and offers a
single-step forecast. It achieves partial proactive capability but
does not interface with any orchestration platform, nor does it address
multi-tenant SLA.

%%─────────────────────────────────────────────────────────────────────
\section{NFV Orchestration-Focused Approaches}
\label{sec:related-nfv}

\textbf{Latency-Aware NFV-6G (MDPI 2024).}
This work focuses on latency-aware VNF placement in 5G/6G networks,
proposing a partial ML model for placement decisions. Orchestration is
a primary contribution, with partial multi-tenant SLA modelling. However,
actual VNF instantiation latency is reported at 8{,}000--14{,}000~ms in
the Kubernetes testbed, and there is no proactive or forecasting component.
\padonap{}'s pre-positioned rate-limiter activates at 505~ms---a factor
of $>15\times$ faster.

\textbf{Cross-Layer SDN/NFV (MDPI 2025).}
This system addresses cross-layer coordination between SDN flow rules
and NFV VNF chains for DDoS response in a DCN evaluation setting.
It achieves full MANO integration and DCN-scale evaluation, but
detection relies on threshold rules rather than ML, and no proactive
or multi-tenant SLA component is described.

%%─────────────────────────────────────────────────────────────────────
\section{Comparison Matrix}
\label{sec:related-matrix}

Table~\ref{tab:related} compares the eight reviewed systems against
\padonap{} on six axes derived directly from the research questions in
Chapter~\ref{ch:intro}: (1) ML-based detection, (2) multi-horizon
forecasting, (3) standards-compliant MANO integration, (4) multi-tenant
SLA guarantee, (5) proactive VNF pre-positioning, and (6) DCN-scale
evaluation.

\begin{table}[H]
\centering
\footnotesize
\caption{Related-work comparison matrix. $\bullet$ = fully supported;
$\circ$ = partially supported; blank = absent.}
\label{tab:related}
\begin{tabular}{@{}lcccccc@{}}
\toprule
System & ML & Multi-horizon & MANO & SLA & Proactive & DCN \\
       & det. & forecast & compliant & guarantee & VNF & eval. \\
\midrule
TST-DDoS (IEEE Access 2025)          & $\bullet$ &            &            &            &            &            \\
NetProbe (ICDCN 2025)                & $\bullet$ &            & $\circ$    &            & $\circ$    & $\bullet$  \\
FlowGuard (NDSS 2022)                & $\bullet$ & $\circ$    &            &            & $\circ$    &            \\
Lightweight ML (Sci.~Rep.~2025)      & $\bullet$ &            &            &            &            &            \\
DAWN (Springer LNCS 2025)            & $\bullet$ &            & $\circ$    &            &            & $\circ$    \\
Transformer+TCN (PLOS One 2025)      & $\bullet$ & $\circ$    &            &            & $\circ$    &            \\
Latency-aware NFV-6G (MDPI 2024)     & $\circ$   &            & $\bullet$  & $\circ$    &            & $\circ$    \\
Cross-layer SDN/NFV (MDPI 2025)      & $\bullet$ &            & $\bullet$  &            &            & $\bullet$  \\
\midrule
\textbf{\padonap{} (ours)}
  & \textbf{$\bullet$}
  & \textbf{$\bullet$}
  & \textbf{$\bullet$}
  & \textbf{$\bullet$}
  & \textbf{$\bullet$}
  & \textbf{$\bullet$} \\
\bottomrule
\end{tabular}
\end{table}

\padonap{} is the only system in this review to achieve a full $\bullet$
on all six axes. The next best-covered systems (NetProbe, DAWN) each
achieve at most three partial or full marks and none covers multi-horizon
forecasting combined with LP-guaranteed SLA.

%%─────────────────────────────────────────────────────────────────────
\section{Research Gaps}
\label{sec:gaps}

The comparison matrix reveals three concrete gaps that this thesis
addresses:

\begin{enumerate}
  \item \textbf{G1 --- No multi-horizon forecast + MANO pairing.}
        No prior system pairs a multi-horizon (4-step) DDoS forecast
        with a standards-compliant (ONAP) MANO lifecycle interface.
        Partial forecasters (FlowGuard, Transformer+TCN) stop at
        single-step and do not interact with real orchestrators.
        \padonap{} closes this gap with Transformer+LSTM at four
        horizons ($t{+}30, 60, 90, 120$~s) integrated into the ONAP
        closed loop.

  \item \textbf{G2 --- No LP-backed SLA under mitigation.}
        None of the reviewed systems formally guarantees per-tenant
        bandwidth floors during active mitigation. This is a critical
        gap in multi-tenant DCN deployments where URLLC tenants must
        retain their contracted floor even when the orchestrator
        activates a VNF chain. \padonap{} closes this gap with an
        LP SLA allocator solved at every tier transition.

  \item \textbf{G3 --- OOD robustness not measured.}
        No prior work explicitly measures the behaviour of their system
        on out-of-distribution traffic (attack classes absent from
        training). This matters because network traffic distributions
        shift over time, and a system that over-escalates on OOD traffic
        causes unnecessary tenant disruption. \padonap{} explicitly
        measures OOD behaviour in scenarios S4 and S5 and demonstrates
        graceful under-response.
\end{enumerate}

%%─────────────────────────────────────────────────────────────────────
\section{Chapter Summary}
\label{sec:related-summary}

This chapter reviewed eight recent DDoS detection and NFV orchestration
systems and identified three research gaps: the absence of a multi-horizon
forecast integrated with MANO, the absence of LP-backed SLA guarantees
under mitigation, and the lack of explicit OOD robustness measurement.
These three gaps directly map to contributions~2, 3, and 4 of \padonap{},
respectively. The proposed system is described in full in the next chapter.
